% BAB IV Hasil dan Pembahasan. Sub-bab baku: Pedoman Tabel 52; isi: Subbab 11.2.2.
\chapter{Hasil dan Pembahasan}
\label{bab:hasil}

\panduan{Susun hasil mengikuti urutan tujuan penelitian, bukan kronologi
pengerjaan. Setiap paragraf hasil memakai pola Assertion-Evidence:
dibuka dengan klaim utuh, diikuti bukti berupa tabel, gambar, persamaan,
atau hasil uji statistik, lalu interpretasi (Subbab 11.2.2). Bab ini
dilarang memuat nama fungsi, sintaks pustaka, atau potongan kode. Kode
hanya ada di lampiran atau repositori (Subbab 11.9 butir 2, temuan
kritis). Bab ini dinilai pada Sub-CPMK3-CPMK5 dengan bobot 26\% dari
Nilai Akhir (Pedoman Tabel 61).}

\section{Deskripsi Data dan Pelaksanaan}
\label{sec:deskripsi-data}

\panduan{Laporkan pelaksanaan pengambilan data: waktu, jumlah unit yang
diperoleh dibandingkan rencana pada BAB~\ref{bab:metode}, data yang
dikeluarkan beserta alasannya, dan hasil uji validitas serta reliabilitas
instrumen.}

Tuliskan ringkasan pelaksanaan pengambilan data di sini.

\section{Hasil per Tujuan Penelitian}
\label{sec:hasil-tujuan}

\panduan{Buat satu subbab untuk setiap tujuan pada
Subbab~\ref{sec:tujuan}. Angka pada teks, tabel, abstrak, dan slide
sidang harus identik (Subbab 13.4.1). Judul kolom memuat satuan, dan
simbol pada tabel sama dengan simbol pada teks (Subbab 11.8). Tabel yang
lebih lebar dari area teks diringkas atau dipindahkan ke lampiran.}

Contoh pola Assertion-Evidence: kelompok responden A memberi skor
kegunaan lebih tinggi daripada kelompok B. Bukti disajikan pada
Tabel~\ref{tab:hasil-sus}, dengan skor dihitung menurut
Persamaan~\eqref{eq:sus}. Data per responden dimuat pada
Lampiran~\ref{lamp:data}.

\begin{table}[htbp]
  \centering
  \caption{Contoh format tabel hasil (angka ilustratif)}
  \label{tab:hasil-sus}
  \begin{tabular}{@{}lrrr@{}}
    \toprule
    Kelompok responden & $n$ & $\bar{S}$ (0--100) & $s$ \\
    \midrule
    Kelompok A & 8 & 72,5 & 9,1 \\
    Kelompok B & 7 & 64,3 & 11,4 \\
    \midrule
    Seluruh responden & 15 & 68,7 & 10,8 \\
    \bottomrule
  \end{tabular}
  \sumber{olahan penulis}
\end{table}

Interpretasikan angka pada tabel tersebut dalam paragraf yang sama.
Jangan biarkan tabel atau gambar berdiri tanpa pembahasan (Subbab 11.8).

\section{Pembahasan dan Perbandingan dengan Literatur}
\label{sec:pembahasan}

\panduan{Bandingkan temuan dengan literatur pada BAB~\ref{bab:tinjauan},
jelaskan penyebab selisih, dan akui keterbatasan (Subbab 11.2.2). Klaim
keunggulan model atau rancangan harus disertai uji signifikansi
statistik beserta nilai $p$ dan ukuran efek (Pedoman Bab 4, temuan
kritis PL1).}

Tuliskan perbandingan temuan Anda dengan penelitian yang dirangkum pada
Tabel~\ref{tab:matriks-sintesis}.

\section{Ancaman Validitas dan Keterbatasan}
\label{sec:ancaman-validitas}

\panduan{Uraikan ancaman validitas internal, eksternal, dan konstruk
beserta langkah mitigasi yang telah dilakukan.}

Tuliskan ancaman validitas dan keterbatasan penelitian di sini.
